%=====================================================
\section{Properties of Relations}
%=====================================================

%-----------------------------------------------------
\begin{frame}{Properties of Relations}

The six fundamental properties are

\begin{itemize}
    \item Reflexive
    \item Irreflexive
    \item Symmetric
    \item Antisymmetric
    \item Asymmetric
    \item Transitive
\end{itemize}

These properties help determine whether a relation is an Equivalence Relation or a Partial Order.

\end{frame}

%-----------------------------------------------------

\begin{frame}{Reflexive and Irreflexive Relations}

\begin{columns}

\column{0.48\textwidth}

\textbf{Reflexive}

\vspace{0.1cm}

$\forall a\in A,\;(a,a)\in R$

Every element is related to itself.

\vspace{0.3cm}

Example:

$\le,\quad =,\quad \subseteq$

\column{0.48\textwidth}

\textbf{Irreflexive}

\vspace{0.1cm}

$\forall a\in A,\;
(a,a)\notin R$

No element is related to itself.

\vspace{0.3cm}

Example:

$<,\quad >$

\end{columns}

\end{frame}

%-----------------------------------------------------

\begin{frame}{Symmetric Relation}

\textbf{Definition}

\[
(a,b)\in R
\Longrightarrow
(b,a)\in R
\]

\vspace{0.5cm}

Example:

\[
xRy
\iff
x
\text{ and }
y
\text{ belong to the same class}
\]

\[
(Student_1,Student_2)
\Longrightarrow
(Student_2,Student_1)
\]

\end{frame}

%-----------------------------------------------------

\begin{frame}{Antisymmetric Relation}

\textbf{Definition}

\[
(a,b)\in R
\land
(b,a)\in R
\Longrightarrow
a=b
\]

\vspace{0.5cm}

Example:
\[
\subseteq
\]
If
\[
A\subseteq B
\; and\;
B\subseteq A
\]
then
\[
A=B
\]

\end{frame}

%-----------------------------------------------------

\begin{frame}{Asymmetric Relation}

\textbf{Definition}

\[
(a,b)\in R
\Longrightarrow
(b,a)\notin R
\]

\vspace{0.4cm}

Example:
\[
<
\]
If
\[
3<5
\]
then
\[
5<3
\]

is false.

\end{frame}

%-----------------------------------------------------

\begin{frame}{Transitive Relation}

\textbf{Definition}

\[
(a,b)\in R
\land
(b,c)\in R
\Longrightarrow
(a,c)\in R
\]

\vspace{0.5cm}

Example:

\[
\le,\quad \subseteq
\]

\end{frame}

%-----------------------------------------------------

\begin{frame}{Applications in AI \& Software Engineering}

\begin{block}{Software Dependency Graph}

\vspace{0.1cm}

If
\[
Module_A
\rightarrow
Module_B
\]
and
\[
Module_B
\rightarrow
Module_C
\]
then
\[
Module_A
\]
indirectly depends on
\[
Module_C
\]
illustrating transitivity.

\end{block}

\end{frame}